%%%%%%%%%%%%%%%%%%%%%%%%%%%%%%%%%%%%%%%%%%%%%%%%%
\section{Quasi-quark PDFs at one-loop order}
\label{sec:oneloop}

In this section, we study the quasi-quark PDFs at one-loop and its renormalization. Feynman diagrams for quasi-quark PDFs of an asymptotic quark of momentum $p$ at one-loop order in the Feynman gauge are shown in Fig.~\ref{fig:oneloop}.  The diagram (a) in Fig.~\ref{fig:oneloop} gives
\begin{align}
\begin{split}\label{eq:1a0}
M_{1a}=&\frac{e^{i p_z \xi_z}}{p_z}\frac{1}{N_c} \text{Tr}_c[T^aT^a] \int_a^{\xi_z-2a}dr_1\int_{r_1+a}^{\xi_z-a}dr_2\\
&\times \int \frac{d^4l}{(2\pi)^4} \,e^{-i p_z \xi_z} \,e^{il_z(r_2-r_1)} \left(\frac{-ig_{\mu\nu}}{l^2}\right) \\
&\times (-ig_s n_z^\mu)(-ig_s n_z^\nu)  \text{Tr}\left[\frac{1}{2}\, \slashed{p}\, \frac{1}{2}\gamma_z\right]   \\
=& \frac{\alpha_s C_F}{4 i \pi^3} \int_a^{\xi_z-2a}dr_1\int_{r_1+a}^{\xi_z-a}dr_2 \int d^4l\frac{e^{il_z(r_2-r_1)}}{l^2}
\end{split}
\end{align}
where we introduced a cutoff ``$a$'' between fields along the gauge link to regularize potential linear UV divergence\footnote{Although explicit result of UV divergence depends on UV regulators, our conclusions, like power counting rules and renormalization structure, are independent of them.}. The upper limit of $r_1$-integration in Eq. (9) is $\xi_z-2a$ because it needs a separation $a$ from the upper limit of $r_2$-integration.  We will show that this cutoff is enough to regularize all UV divergence in this diagram, and thus we do not introduce DR here. For the following calculation, we introduce a vector $\bar{l}^\mu$, which is the same as $l^\mu$ except that it does not have the $z$-direction component: $l^\mu = \bar{l}^\mu + l_z n_z^\mu$ and $l^2=\bar{l}^2-l_z^2$. We will refer the integration of $\bar{l}^\mu$ as three-dimensional (3-D) integration and the integration of $l^\mu$ as four-dimensional (4-D) integration.  With $d^4l = d^{3}\bar{l}\, dl_z$, let's consider the following phase space integration,
\begin{eqnarray}\label{eq:3dexpansion}
\int \frac{d^{3}\bar{l}}{l^2} &=&
\int \frac{d^{3}\bar{l}}{\bar{l}^2-l_z^2}
\nonumber\\
&=& \int d^{3}\bar{l} \left( \frac{1}{\bar{l}^2}+ \frac{l_z^2}{(\bar{l}^2-l_z^2)\bar{l}^2}\right).
\end{eqnarray}
The first term in Eq.~(\ref{eq:3dexpansion}) is linear divergent, but its coefficient is proportional to $\int dl_z e^{il_z(r_2-r_1)} = 2\pi \delta(r_2-r_1)$, which vanishes because $r_2$ and $r_1$ cannot be at the same point. Thus, effectively, the 3-D integration above is finite. On the other hand, it is easy to see that, if keeping $\bar{l}^\mu$ finite, the integration of $l_z$ is also finite even if $a\to 0$. Therefore, the Eq.~\eqref{eq:1a0} can be UV divergent only in the region where all components of $l^\mu$ go to infinity. As a result, the spacing ``$a$'' can regularize all UV divergences, which gives
\begin{align}\label{eq:1a}
\begin{split}
M_{1a}\overset{\text{div}}{=}-\frac{\alpha_s C_F}{\pi}\frac{|\xi_z|}{a}+\frac{\alpha_s C_F}{\pi}\ln \frac{|\xi_z|}{a},
\end{split}
\end{align}
where we have included the situation when $\xi_z < 0$.


%%%%%%%%%%%%%%%%%%%%%%%%%%%%%%%
\begin{figure}[t]
\begin{center}
\includegraphics[width=2.5in]{oneloop}
\caption{\label{fig:oneloop}
Feynman diagrams for quasi-quark PDFs of an asymptotic quark of momentum $p$ at one-loop order.}
\end{center}
\end{figure}
%%%%%%%%%%%%%%%%%%%%%%%%%%%%%%


From Fig.~\ref{fig:oneloop}, we have the contribution from the diagram (b) as
\begin{align}
\begin{split}\label{eq:1b0}
M_{1b}=& g_s^2C_F \int_a^{\xi_z-a}dr\int \frac{d^4l}{(2\pi)^4}\,
\frac{e^{il_z r}}{l^2(p-l)^2}
 \\
& \times \frac{1}{p_z} \text{Tr}\left[\frac{1}{2}\, \slashed{p}\, \frac{1}{2}\gamma_z(\slashed{p}-\slashed{l})\gamma_z\right],
\end{split}
\end{align}
where we again use the spacing ``$a$'' as a regulator. For the 3-D integration, it has potential logarithmic divergence from the term proportional to $\slashed{l}$,
\begin{align}
\int \frac{l^\mu\,d^{3}\bar{l}}{l^2(p-l)^2}.
\end{align}
However, the above integration is finite because, to extract the potential logarithmic divergence, we can set the external physical scale $p\to 0$ in the above integration and thus it becomes an odd function in $\bar{l}^\mu$, which vanishes under integration. This explains why the spacing $a$ can regularize the UV divergence of Eq.~\eqref{eq:1b0}. The UV divergence of the diagram (b) is
\begin{align}\label{eq:1b}
\begin{split}
M_{1b}\overset{\text{div}}{=}-\frac{\alpha_s C_F}{2\pi} \ln \frac{|\xi_z|}{a},
\end{split}
\end{align}
where we again included the situation when $\xi_z<0$.

Similarly, we have from the diagram (c) in Fig.~\ref{fig:oneloop},
\begin{align}\label{eq:1c0}
\begin{split}
M_{1c}=&\frac{i g_s^2 \mu_r^{2\epsilon}}{2p_z}(1-\epsilon)C_F\int\frac{d^dl}{(2\pi)^d} \frac{\text{Tr}[\slashed{p}\gamma_z {\slashed{p}}(\slashed{p}-\slashed{l})]}{p^2l^2(p-l)^2}\\
=&i g_s^2 \mu_r^{2\epsilon} C_F(1-\epsilon)\int\frac{d^dl}{(2\pi)^d} \frac{1}{l^2(p-l)^2},
\end{split}
\end{align}
where space-time dimension is defined as $d=4-2\epsilon$. It is clear that the integral in Eq.~(\ref{eq:1c0}) vanishes in DR when $p^2=0$, due to the apparent cancelation between the UV and IR poles in $1/\epsilon$.  The UV divergence from this diagram is logarithmic only if all components of $l^\mu$ go to infinity, and is given in DR by
\begin{align}\label{eq:1c}
\begin{split}
M_{1c}\overset{\text{div}}{=}-\frac{\alpha_s C_F}{4\pi} \frac{1}{\epsilon}.
\end{split}
\end{align}
It is well-known that, at this order and higher orders, UV divergence of massless quark self-energy diagrams can be removed by the renormalization of quark field.

The diagram (d) in Fig.~\ref{fig:oneloop} gives
\begin{align}
\begin{split}\label{eq:1d0}
M_{1d}=& -i g_s^2C_F e^{i p_z\xi_z} \int \frac{d^4k}{(2\pi)^4} e^{-ik_z \xi_z} \frac{1}{k^4 (p-k)^2} \\
&\times \frac{1}{4p_z } \text{Tr}[\slashed{p}\gamma_\mu\slashed{k}\gamma_z\slashed{k}\gamma^\mu]
\end{split}
\end{align}
where $k=p-l$.
Because of the oscillation factor $e^{-ik_z \xi_z}$, contribution from the region where $k_z$ goes to infinity is highly suppressed. Thus, it is UV finite at finite $\xi_z$, although it is divergent as $\xi_z \to 0$,
\begin{align}\label{eq:1d}
\begin{split}
M_{1d}\overset{\text{div}}{=}-\frac{\alpha_s C_F}{2\pi} \ln(|\xi_zp_z|).
\end{split}
\end{align}

In summary, the total one-loop contribution to the UV divergence of quasi-quark PDFs of a quark of momentum $p$ at an arbitrary $\xi_z$ is
\begin{align}\label{eq:1}
\begin{split}
M^{(1)}&\overset{\text{div}}{=}M_{1a}+2\times M_{1b}+ 2\times\frac{1}{2}M_{1c}+M_{1d}\\
&=\frac{\alpha_s C_F}{\pi}\left(-\frac{|\xi_z|}{a} - \frac{1}{4\epsilon} - \frac{1}{2} \ln(|\xi_zp_z|) \right),
\end{split}
\end{align}
where the $ \ln(|\xi_zp_z|)$ term is not UV divergent, but divergent as $\xi_z \to 0$
\footnote{If we take $\xi_z\to 0$, what we calculated is an one-loop correction to a local vector current.  The finite $\xi_z$ in Eq.~\eqref{eq:1} effectively regularizes the UV divergence of the one-loop vertex diagram  Fig. \ref{fig:oneloop}(d), while UV divergence of the self-energy diagram in Fig. \ref{fig:oneloop}(c) is regularized by the dimensional regularization.  From Fig. \ref{fig:oneloop}(d), we obtained the $\ln(|\xi_z p_z|)$ term after we took $\epsilon\to 0$ with $\xi_z$ fixed. If we need to take $\xi_z\to 0$ to mimic the local current, we should keep $\epsilon$ finite to regularize the one-loop UV divergence, which changes $\ln(|\xi_z p_z|)$ to $ - \frac{1}{2\epsilon}$.  As a result, we find from Eq.~\eqref{eq:1} that UV divergences of the vector current vanish at one-loop level if we take $\xi_z\to 0$ with fixed $\epsilon$ and $a$, which is consistent with the expectation of current conservation.}.
 We find that, at this order, UV divergences only come from the region where all loop momenta go to infinity, thus UV divergences are localized in coordinate space.  We will show in the next section that this behavior remains true up to all orders in QCD perturbation theory.


%%%%%%%%%%%%%%%%%%%%%%%%%%%%%%%
\begin{figure}[ht]
\begin{center}
\includegraphics[width=3.3in]{oneloopg}
\caption{\label{fig:oneloopg}
Feynman diagram for quasi-quark PDFs of an asymptotic gluon of momentum $p$ at one-loop order.}
\end{center}
\end{figure}
%%%%%%%%%%%%%%%%%%%%%%%%%%%%%%

Feynman diagrams for quasi-quark PDFs of an asymptotic gluon of momentum $p$ at one-loop order in the Feynman gauge are shown in Fig.~\ref{fig:oneloopg}, plus the complex conjugate diagram of (b). A general argument in the next section will show that UV divergences of all diagrams of quasi-PDFs can only come from the 4-D integration, and thus localized in coordinate space. If we keep $\xi_z$ to be finite, all one-loop diagrams in Fig.~\ref{fig:oneloopg} cannot be local, just like the diagram (d) in Fig.~\ref{fig:oneloop}, and therefore, all these one-loop diagrams in Fig.~\ref{fig:oneloopg} must be UV-finite.  But, they can be divergent as $\xi_z\to 0$. To demonstrate this feature, let's take the diagram (a) in Fig.~\ref{fig:oneloopg} as an example. The Fig.~\ref{fig:oneloopg}(a) gives
\begin{align}
\begin{split}\label{eq:2a0}
M_{2a}\propto& \int_0^{\xi_z} dr_1 \int_{r_1}^{\xi_z} dr_2 \int d^4l\, e^{-il_z \xi_z} \frac{l_z}{l^2} \\
=&\frac{\xi_z^2}{2}\int dl_z \, e^{-il_z \xi_z} \, l_z \int d^{3}\bar{l} \left( \frac{1}{\bar{l}^2}+ \frac{l_z^2}{(\bar{l}^2-l_z^2)\bar{l}^2}\right),
\end{split}
\end{align}
which seems to be linearly UV divergent for the 3-D integration. However, similar to the case of diagram (a) in Fig.~\ref{fig:oneloop}, the linear divergent term is proportional to $\delta^\prime(\xi_z)$, which vanishes for finite $\xi_z$. The second term is finite under integration of $\bar{l}$, which gives
\begin{align}
\begin{split}\label{eq:2a}
& \frac{\xi_z^2}{2}\int dl_z \, e^{-il_z \xi_z} \, l_z \int d^{3}\bar{l} \frac{l_z^2}{(\bar{l}^2-l_z^2)\bar{l}^2}\\
\propto&\frac{\xi_z^2}{2}\int dl_z \, e^{-il_z \xi_z} \, \frac{l_z^3}{|l_z|}\\
=&\frac{2i}{\xi_z},
\end{split}
\end{align}
where the proportional relation can be simply obtained by dimensional counting. We therefore find that Fig.~\ref{fig:oneloopg}(a) is UV finite, but behaves as $1/\xi_z$ when $\xi_z \to 0$. Similarly, we found that Figs.~\ref{fig:oneloopg}(b,c) are also UV finite, which indicates that at one-loop, the quasi-quark PDFs get no mixing from a gluon under the UV renormalization.


Before continue, we want to note that the UV behavior found above is significantly different from that of PDFs.  UV divergences of PDFs come from the region of loop momentum $(l_+, l_-, \vec{l}_\perp)\sim(1,\lambda^2,\lambda)$ when $\lambda\to\infty$. Thus, with the momentum component $l_+\sim {\cal O}(1)$, the UV divergence for the normal PDFs is nonlocal in the ``$-$'' direction in coordinate space. It is this fact that the renormalization of PDFs is a convolution, mixed with all twist-2 PDFs with operators along the ``$-$'' direction, while the renormalizaiton of quasi-PDFs is a multiplicative factor as we will show.

We further note that, even if UV divergences ($1/a$, $\ln(a)$ or $1/\epsilon$) are renormalized, the $\ln(|\xi_z|)$ dependence (see Eq.~\ref{eq:1d} as an example) and $1/\xi_z$ dependence (see Eq.~\ref{eq:2a} as an example) of one-loop diagrams signal the divergence of quasi-PDFs as $\xi_z\to 0$, which causes the difficulty in getting a well-defined momentum-space quasi-PDFs if one simply Fourier transforms the coordinate-space quasi-PDFs, as mentioned in the last section. This behaviour will certainly remain true to higher orders in QCD perturbation theory.
